% Edmonds--Karp 主要思想（Python 验算）：外弧 s->a,s->b,a->t,b->t 容量 4，中间弧 a->b 容量 1。
% (a) 交替选 s->a->b->t 与 s->b->a->t（经反向残量 b->a），每次只送 1，流值 8 需 8 次增广。
% (b) BFS 层 h(s)=0, h(a)=h(b)=1, h(t)=2；两条两弧路径各送 4，两次增广达到流值 8。
\begin{tikzpicture}[>=Stealth,
  v/.style={circle,draw=FigureStroke,fill=FigureBlueFill,line width=.8pt,minimum size=6mm,inner sep=0pt},
  arc/.style={->,draw=FigureGrid,line width=.8pt},
  p1/.style={->,draw=FigureBlue,line width=1.6pt},
  p2/.style={->,draw=FigureGreen,dashed,line width=1.2pt},
  aug/.style={->,draw=FigureYellow,line width=1.8pt},
  w/.style={fill=white,inner sep=.5pt,font=\figurefont\fontsize{7}{8}\selectfont},
  lab/.style={font=\figurefont\fontsize{7.5}{9}\selectfont},
  lev/.style={font=\figurefont\fontsize{7}{8}\selectfont,text=FigureMuted}]
\newcommand{\verts}{%
  \node[v] (s) at (0,0) {$s$}; \node[v] (a) at (1.4,.95) {$a$};
  \node[v] (b) at (1.4,-.95) {$b$}; \node[v] (t) at (2.8,0) {$t$};}
\begin{scope}
  \verts
  \draw[p1] (s) to[bend left=10] node[w]{$4$}(a);
  \draw[p1] (a) to[bend right=14] node[w,pos=.35]{$1$}(b);
  \draw[p1] (b) to[bend right=10] node[w]{$4$}(t);
  \draw[p2] (s) to[bend right=10] node[w]{$4$}(b);
  \draw[p2] (b) to[bend right=14] node[w,pos=.35]{$1$}(a);
  \draw[p2] (a) to[bend left=10] node[w]{$4$}(t);
  \node[lab,align=center] at (1.4,-1.85) {(a) 交替沿两条三弧路\\每次只送$1$};
\end{scope}
\begin{scope}[xshift=4.2cm]
  \verts
  \draw[aug] (s)--node[w]{$4$}(a); \draw[aug] (a)--node[w]{$4$}(t);
  \draw[aug] (s)--node[w]{$4$}(b); \draw[aug] (b)--node[w]{$4$}(t);
  \draw[arc] (a)--node[w]{$1$}(b);
  \foreach \x/\h in {0/0,1.4/1,2.8/2} \node[lev] at (\x,1.55) {$h=\h$};
  \node[lab,align=center] at (1.4,-1.85) {(b) 先找两弧路\\两次各送$4$};
\end{scope}
\end{tikzpicture}
